% Figure — Bandgap de Brokaw de l'exercice « Une référence des circuits intégrés », découpé en quadripôles A et β_F (EANA3, chapitre 3).
% D'après le sujet [td03-2026, p. 1, figure 1], redessiné.
% Rendu SVG :
%   latex -output-directory=<dossier temporaire> 03-bandgap-brokaw.tex
%   dvisvgm --no-fonts --exact-bbox -o 03-bandgap-brokaw.svg <dossier temporaire>/03-bandgap-brokaw.dvi
\documentclass[border=4pt]{standalone}
\usepackage[T1]{fontenc}
\usepackage[utf8]{inputenc}
\usepackage{amsmath}
\usepackage{sansmath}
\renewcommand{\familydefault}{\sfdefault}
\usepackage[european, straightvoltages, americancurrents]{circuitikz}
\usetikzlibrary{shapes.geometric,calc}
\begin{document}
\sansmath
\begin{circuitikz}[font=\small, >=latex, fleche/.style={->, thick}]
  % bloc A
  \draw[dashed, thick] (1.6,3.0) rectangle (7.6,6.4);  \node[font=\bfseries] at (7.3,6.15) {$A$};
  \draw (2.8,4.6) coordinate (np) to[R, l=$R_c$] (2.8,5.9) -- (3.6,5.9) node[vcc]{$V_{cc}$} -- (4.4,5.9) to[R, l=$R_c$] (4.4,3.8) coordinate (nm);
  \node[draw, thick, isosceles triangle, minimum height=12mm, fill=black!4, inner sep=1pt] (Av) at (5.9,4.2) {$A_v$};
  \draw (np) -- (4.28,4.6) arc[start angle=180, end angle=0, radius=0.12] -- (Av.west |- np);  \draw (nm) -- (Av.west |- nm);
  \node[font=\scriptsize] at (5.3,4.75) {$+$};  \node[font=\scriptsize] at (5.3,3.65) {$-$};
  \draw[<->] (5.05,3.85) -- node[left, font=\footnotesize]{$\varepsilon$} (5.05,4.55);
  \draw (Av.east) -- (8.6,4.2) coordinate (vr) -- (9.4,4.2) node[ocirc]{};
  \fill (vr) circle (1.5pt);
  % entrées
  \draw (0,4.6) -- (np);  \draw (0,3.8) -- (nm);
  \draw[fleche, blue!70!black] (0.0,4.85) -- node[above]{$i_i$} (0.7,4.85);
  \node[above] at (2.0,2.4) {$i_1$};  \node[above] at (2.0,1.8) {$i_2$};
  % bloc beta
  \draw[dashed, thick] (1.6,-2.8) rectangle (7.6,2.8);  \node[font=\bfseries] at (7.2,2.55) {$\beta_F$};
  \draw (4.0,0.6) node[npn, anchor=B, xscale=-1] (Q2) {};  \node[left] at (2.75,1.0) {$Q_2$ ($\times 8$)};
  \draw (5.0,0.6) node[npn, anchor=B] (Q1) {};  \node[right] at (6.0,0.2) {$Q_1$};
  \draw (1.0,4.6) node[circ]{} -- (1.0,2.4) -| (Q1.C);  \draw (0.6,3.8) node[circ]{} -- (0.6,1.8) -| (Q2.C);
  \draw (Q2.B) -- (Q1.B);  \draw (4.5,0.6) node[circ]{} -- (4.5,1.5) -- ($(Q1.C |- 0,1.5) + (-0.12,0)$) arc[start angle=180, end angle=0, radius=0.12] -- (8.6,1.5) -- (vr);
  \draw (Q2.E) to[R, l=$R_2$] (Q2.E |- 0,-1.4) coordinate (e) -- (Q1.E |- e) -- (Q1.E);
  \draw[fleche] (2.5,-1.3) -- node[left]{$\Delta v_{be}$} (2.5,-0.2);
  \draw (4.5,-1.4) node[circ]{} to[R, l=$R_1$] (4.5,-3.4);
  \draw (4.5,-3.4) node[ground]{} -- (9.4,-3.4) node[ocirc]{};
  \draw[fleche] (9.7,-3.2) -- node[right]{$v_{ref}$} (9.7,4.0);
\end{circuitikz}
\end{document}
